\cleardoublepage
\chapter{STUDI LITERATUR}
\label{chap:studi-literatur}

\section{Standar Regulasi Bangunan Gedung Cerdas (BGC)}
Penyelenggaraan Bangunan Gedung Cerdas (BGC) di Indonesia secara hukum diatur dalam Peraturan Menteri Pekerjaan Umum dan Perumahan Rakyat Nomor 10 Tahun 2023 \autocite{PermenPUPR_10_2023}. Regulasi tersebut menetapkan sembilan prinsip utama BGC, antara lain beroperasi secara otomatis, saling terhubung dan terintegrasi, serta memiliki perlindungan yang kuat terhadap ancaman siber. Untuk tahap pemanfaatan bangunan, Surat Edaran Menteri Pekerjaan Umum Nomor 22/SE/M/2024 \autocite{pu2024se22} memberikan pedoman teknis yang lebih spesifik melalui enam parameter penilaian kinerja:

\begin{enumerate}
    \item Keamanan Siber: Menilai keandalan sistem informasi dalam menangani ancaman digital, kebocoran data, dan manajemen otentikasi.
    \item Protokol dan Jaringan Komunikasi: Menilai standarisasi arsitektur jaringan lokal untuk mendukung interoperabilitas antarperangkat keras.
    \item Integrasi Data dan Sistem: Menilai kemampuan sistem penengah dalam mengonsolidasikan berbagai subsistem tertutup menjadi satu kesatuan pusat data.
    \item Manajemen Operasi: Menilai efisiensi pengelolaan fasilitas operasional secara terpusat.
    \item Dampak terhadap Manusia: Menilai kontribusi teknologi terhadap kenyamanan, privasi, dan keselamatan pengguna bangunan.
    \item Kemampuan Sistem: Menilai fungsionalitas dari 16 elemen teknis BGC. Sistem kontrol akses (\textit{access control system}) dan sistem kamera pengawas (\textit{closed circuit television}) merupakan bagian dari komponen vital kemampuan sistem ini.
\end{enumerate}

\section{Sistem Kontrol Akses dan Interoperabilitas IoT}
Sistem kontrol akses fisik berfungsi sebagai baris pertahanan pertama dalam manajemen keamanan bangunan modern \autocite{norman2011electronic}. Penggunaan teknologi biometrik seperti pengenalan wajah (\textit{face recognition}) menawarkan tingkat validitas yang jauh lebih tinggi dibandingkan kredensial konvensional seperti kartu RFID yang rentan dipindahtangankan \autocite{wanjare2024}.

Namun, tantangan terbesar pada implementasi perangkat keras skala industri adalah masalah \textit{vendor lock-in} dan keterisolasian data (\textit{data silo}) \autocite{s19020229}. Perangkat lunak bawaan pabrikan umumnya berbasis desktop lokal dan menggunakan ekosistem tertutup. Kondisi tersebut menghambat pertukaran informasi secara terbuka, sehingga log pergerakan personel tidak dapat ditarik secara seketika (\textit{real-time}) oleh platform luar seperti \textit{Building Management System} (BMS). Oleh karena itu, diperlukan arsitektur lapisan penengah (\textit{middleware}) berbasis standar terbuka untuk menjembatani ekosistem perangkat keras tertutup dengan aplikasi web modern.

\section{Protokol \textit{Intelligent Security API} (ISAPI)}
\textit{Intelligent Security API} (ISAPI) adalah protokol komunikasi berbasis standar HTTP/REST yang dikembangkan secara spesifik untuk integrasi perangkat keras keamanan \autocite{hikvision_isapi_person}. Berbeda dengan \textit{Software Development Kit} (SDK) konvensional yang mengikat aplikasi pada dependensi pustaka (\textit{library}) dan sistem operasi tertentu, ISAPI bersifat mandiri platform (\textit{platform-independent}). Protokol ini memungkinkan peladen melakukan kontrol penuh terhadap memori dan sensor terminal melalui pemanggilan rute titik akhir (\textit{endpoint outbound}) menggunakan format data terstruktur seperti JSON atau XML.

Keamanan transmisi data pada protokol ISAPI diperkuat melalui mekanisme \textit{Digest Authentication}. Metode otentikasi ini menangani pertukaran kredensial secara aman melalui skema tantangan kriptografi (\textit{cryptographic challenge-response}). Saat peladen mengirimkan permintaan awal, terminal akan menolak secara terprogram dengan memberikan respons HTTP 401 beserta parameter sesi dinamis berupa \textit{nonce}, \textit{realm}, dan \textit{qop}. Peladen kemudian wajib menghitung ulang nilai \textit{hash} MD5 yang menggabungkan kata sandi perangkat dengan parameter tantangan tersebut sebelum melakukan pengiriman ulang (\textit{retry}). Mekanisme ini memastikan kredensial asli tidak pernah ditransmisikan dalam bentuk teks terang (\textit{plain text}), sehingga memitigasi risiko serangan penyadapan (\textit{man-in-the-middle attack}).

\section{Arsitektur Perangkat Lunak \textit{Middleware} dan Komunikasi Web}
Pengembangan sistem manajemen terpusat berskala industri menuntut pemisahan logika program yang disiplin guna menjaga skalabilitas.

\subsection{NestJS sebagai Lapisan \textit{Middleware}}
NestJS merupakan kerangka kerja (\textit{framework}) \textit{backend} Node.js yang mengadopsi arsitektur modular yang kaku berbasis Pemrograman Berorientasi Objek (OOP) dan \textit{Dependency Injection} \autocite{nestjs_docs}. Berjalan di atas mesin \textit{runtime} V8, kerangka kerja ini sangat efisien dalam menangani operasi masukan/keluaran yang bersifat \textit{non-blocking} (\textit{asynchronous non-blocking I/O}). Karakteristik tersebut membuatnya ideal untuk mengorkestrasi tugas latar belakang (\textit{background process} atau \textit{cron job}) seperti penarikan log berkala dari banyak terminal gerbang secara simultan.

\subsection{Next.js sebagai Lapisan Presentasi}
Next.js merupakan kerangka kerja tingkat lanjut berbasis React yang mendukung arsitektur \textit{Server-Side Rendering} (SSR) \autocite{nextjs_docs}. Melalui SSR, proses perendiran komponen antarmuka dilakukan secara langsung di sisi peladen sebelum dikirimkan ke peramban klien. Mekanisme ini menjamin performa pemuatan \textit{dashboard} manajemen tetap responsif meskipun sistem harus menyajikan visualisasi data log dalam volume besar.

\subsection{Komunikasi Dua Arah dengan WebSocket dan Socket.IO}
Protokol HTTP konvensional memiliki keterbatasan dalam menyajikan data seketika karena bersifat satu arah (\textit{request-response}). Untuk memfasilitasi pemantauan instan tanpa pemuatan ulang halaman (\textit{refresh}), teknologi WebSocket digunakan untuk menyediakan saluran komunikasi \textit{full-duplex} yang persisten. Socket.IO bertindak sebagai lapisan abstraksi di atas WebSocket yang menyediakan fitur manajemen jabat tangan (\textit{handshake}) terotentikasi, penanganan putus koneksi otomatis, serta diseminasi data secara instan (\textit{broadcast}) dari peladen ke banyak pengelola gedung \autocite{socketio_docs}.

\section{Lapisan Keamanan Jaringan dan Manajemen Basis Data}

\subsection{NGINX sebagai \textit{Reverse Proxy}}
NGINX bertindak sebagai lapisan pelindung terluar infrastruktur (\textit{edge security}) \autocite{nginx_docs}. Sebagai \textit{reverse proxy}, NGINX bertindak sebagai perantara tunggal yang menerima seluruh permintaan dari luar jaringan lokal sebelum diteruskan ke peladen aplikasi NestJS. Konfigurasi ini menyembunyikan topologi jaringan internal, mengisolasi peladen utama dari paparan publik langsung, serta dapat dikonfigurasi untuk membatasi ukuran muatan (\textit{payload size}) guna memitigasi ancaman serangan kelumpuhan jaringan (\textit{Denial of Service}).

\subsection{Integritas Relasional PostgreSQL dan Prisma ORM}
Manajemen data log pergerakan manusia menuntut konsistensi data yang mutlak. PostgreSQL sebagai Sistem Manajemen Basis Data Relasional (RDBMS) menjamin kepatuhan penuh terhadap sifat ACID (\textit{Atomicity, Consistency, Isolation, Durability}) \autocite{postgresql_docs}. Penggunaan \textit{Object-Relational Mapping} (ORM) seperti Prisma mempercepat proses manipulasi data melalui pemodelan skema yang bertipe aman (\textit{type-safe}) \autocite{prisma_docs}. Penegakan aturan integritas referensial seperti \textit{Restricted Delete} memastikan data log historis pada tabel transaksi tidak akan terhapus atau cacat meskipun profil induk pengguna pada tabel relasi mengalami perubahan administratif.
